\subsection{\textsf{p6m}}\label{subsec:p6m}

We let $\Pi = \Z a \oplus \Z b \subset \R^2$ be the lattice spanned by
$a = \left( \begin{matrix} 1 \\ 0 \end{matrix} \right)$ and
$b = \left( \begin{matrix} 1/2 \\ \sqrt{3}/2 \end{matrix} \right)$.
The point group $P$ is
$D_6 =\langle C, \sigma_1 \,|\, C^6, \sigma_1^2, \sigma_1C\sigma_1C \rangle=
\{1, C, C^2, C^3, C^4, C^5,\sigma_1, \sigma_2, \sigma_3, \sigma_4, \sigma_5, \sigma_6\}$, where $\sigma_\ell = C^{\ell - 1}\sigma_1$. This group acts on $\Pi$ and $\R^2$ through the inclusion $D_6 \subset {\rm O}(2)$ def\/ined by
\begin{gather*}
C = \left(
\begin{matrix}
1/2 & -\sqrt{3}/2 \\
\sqrt{3}/2 & 1/2
\end{matrix}
\right), \qquad \sigma_1 =
\left(
\begin{matrix}
1 & 0 \\
0 & -1
\end{matrix}
\right).
\end{gather*}
If we use the identif\/ications $a = 1$ and $b = \tau = \exp 2\pi i/6$ under $\R^2 = \C$, then the actions of $C \in D_6$ and $\sigma_1$ are given by the multiplication by $\tau$ and the complex conjugation, respectively. The closure of a fundamental domain is $\{ s a + t b \,|\, 0 \le s, t \le 1 \}$ or equivalently $\{ s + t \tau \,|\, 0 \le s, t \le 1 \}$. We decompose this region to def\/ine a $D_6$-CW decomposition of $T^2$ as follows:
$$
\begin{array}{|c|c|c|}
\hline
\mbox{0-cell} & \mbox{1-cell} & \mbox{2-cell} \\
\hline
\tilde{e}^0_0 = \pt &
\tilde{e}^1_{01} = (D_6/\{ 1, \sigma_1 \}) \times e^1 &
\tilde{e}^2 = D_6 \times e^2 \tsep{2pt}\\
\tilde{e}^0_1 = D^6/\{ 1, C^3, \sigma_1, \sigma_4 \} &
\tilde{e}^1_{02} = (D_6/\{ 1, \sigma_2 \}) \times e^1 & \\
\tilde{e}^0_2 = D^6/\{ 1, C^2, C^4, \sigma_2, \sigma_4, \sigma_6 \} &
\tilde{e}^1_{12} = (D_6/\{ 1, \sigma_4 \}) \times e^1 & \bsep{2pt}\\
\hline
\end{array}
$$

\begin{itemize}\itemsep=0pt
\item
($0$-cell)
The $0$-cell $\tilde{e}^0_0 = (D_6/D_6) \times e^0 = \pt$ is the unique f\/ixed point on $T^2$. The other $0$-cells are def\/ined as follows:
\begin{gather*}
\tilde{e}^0_1 =\left\{\frac{1}{2}, \frac{\tau}{2}, \frac{1+\tau}{2}\right\}\cong \big(D_6/\big\{ 1, C^3, \sigma_1, \sigma_4 \big\}\big) \times e^0, \\
\tilde{e}^0_2 =\left\{\frac{1 + \tau}{3}, \frac{2(1+\tau)}{3}\right\}\cong \big(D_6/\big\{ 1, C^2, C^4, \sigma_2, \sigma_4, \sigma_6 \big\}\big) \times e^0.
\\
\xygraph{
!{<0cm,0cm>;<1cm,0cm>:<0cm,1cm>::}
!{(0,0)}*+{\bullet}="a"
!{(2,0)}="b"
!{(1,1.7)}="c"
!{(3,1.7)}="d"
!{(1,0.6)}="e"
!{(2,1.2)}="f"
"a"-@{.}"b"
"a"-@{.}"c"
"b"-@{.}"c"
"b"-@{.}"d"
"c"-@{.}"d"
}
\qquad
\xygraph{
!{<0cm,0cm>;<1cm,0cm>:<0cm,1cm>::}
!{(0,0)}="a"
!{(2,0)}="b"
!{(1,1.7)}="c"
!{(3,1.7)}="d"
%!{(1,0.6)}="e"
%!{(2,1.2)}="f"
!{(1,0)}*+{\bullet}="g"
!{(0.5,0.85)}*+{\bullet}="h"
%!{(2,1.7)}="i"
%!{(2.5,0.85)}="j"
!{(1.5,0.85)}*+{\bullet}="k"
"a"-@{.}"g"
"g"-@{.}"b"
"a"-@{.}"h"
"h"-@{.}"c"
"b"-@{.}"k"
"k"-@{.}"c"
"b"-@{.}"d"
"c"-@{.}"d"
}
\qquad
\xygraph{
!{<0cm,0cm>;<1cm,0cm>:<0cm,1cm>::}
!{(0,0)}="a"
!{(2,0)}="b"
!{(1,1.7)}="c"
!{(3,1.7)}="d"
!{(1,0.6)}*+{\bullet}="e"
!{(2,1.2)}*+{\bullet}="f"
"a"-@{.}"b"
"a"-@{.}"c"
"b"-@{.}"c"
"b"-@{.}"d"
"c"-@{.}"d"
}
\end{gather*}

\item
($1$-cell)
For $0 \le i < j \le 2$, the $1$-cell $\tilde{e}^1_{ij}$ consists of the six segments connecting $\tilde{e}^0_i$ and $\tilde{e}^0_j$. They are of the forms $\tilde{e}^1_{01} = (D_6/\{ 1, \sigma_1 \}) \times e^1$, $\tilde{e}^1_{02} = (D_6/\{ 1, \sigma_2 \}) \times e^1$ and $\tilde{e}^1_{12} = (D_6/\{ 1, \sigma_4 \}) \times e^1$.
\begin{gather*}
\xygraph{
!{<0cm,0cm>;<1cm,0cm>:<0cm,1cm>::}
!{(0,0)}*+{\circ}="a"
!{(2,0)}*+{\circ}="b"
!{(1,1.7)}*+{\circ}="c"
!{(3,1.7)}="d"
%!{(1,0.6)}="e"
%!{(2,1.2)}="f"
!{(1,0)}*+{\circ}="g"
!{(0.5,0.85)}*+{\circ}="h"
%!{(2,1.7)}="i"
%!{(2.5,0.85)}="j"
!{(1.5,0.85)}*+{\circ}="k"
"a"-"g"
"g"-"b"
"a"-"h"
"h"-"c"
"b"-"k"
"k"-"c"
"b"-@{.}"d"
"c"-@{.}"d"
}
\qquad
\xygraph{
!{<0cm,0cm>;<1cm,0cm>:<0cm,1cm>::}
!{(0,0)}*+{\circ}="a"
!{(2,0)}*+{\circ}="b"
!{(1,1.7)}*+{\circ}="c"
!{(3,1.7)}*+{\circ}="d"
!{(1,0.6)}*+{\circ}="e"
!{(2,1.2)}*+{\circ}="f"
"a"-@{.}"b"
"a"-@{.}"c"
"b"-@{.}"c"
"b"-@{.}"d"
"c"-@{.}"d"
"a"-"e"
"b"-"e"
"c"-"e"
"d"-"f"
"b"-"f"
"c"-"f"
}
\qquad
\xygraph{
!{<0cm,0cm>;<1cm,0cm>:<0cm,1cm>::}
!{(0,0)}="a"
!{(2,0)}="b"
!{(1,1.7)}="c"
!{(3,1.7)}="d"
!{(1,0.6)}*{\circ}="e"
!{(2,1.2)}*{\circ}="f"
!{(1,0)}*+{\circ}="g"
!{(0.5,0.85)}*+{\circ}="h"
!{(2,1.7)}*+{\circ}="i"
!{(2.5,0.85)}*+{\circ}="j"
!{(1.5,0.85)}*+{\circ}="k"
"e"-"g"
"e"-"h"
"e"-"k"
"f"-"k"
"f"-"i"
"f"-"j"
"a"-@{.}"g"
"g"-@{.}"b"
"a"-@{.}"h"
"h"-@{.}"c"
"b"-@{.}"k"
"k"-@{.}"c"
"b"-@{.}"j"
"j"-@{.}"d"
"c"-@{.}"i"
"i"-@{.}"d"
}
\end{gather*}

\item
($2$-cell)
The $2$-cell $\tilde{e}^2 = D_6 \times e^2$ consists of the twelve small triangular regions surrounded by the $1$-cells.
\end{itemize}

Let $Y \subset T^2$ be the invariant subspace $Y = \tilde{e}^0_0 \cup \tilde{e}^0_1 \cup \tilde{e}^1_{01}$.


\begin{lem} \label{lem:p6m_Y}
The equivariant cohomology of $Y$ is given as follows:
$$
\begin{array}{|c|c|c|c|c|}
\hline
 & n = 0 & n = 1 & n = 2 & n = 3 \\
\hline
H^n_{D_6}(Y; \Z) & \Z & 0 & \Z_2^{\oplus 3} & \Z_2^{\oplus 2} \tsep{2pt}\bsep{2pt} \\
\hline
\end{array}
$$
\end{lem}

\begin{proof}
We can f\/ind $D_6$-invariant subspaces $U$ and $V$ in $Y$ which have the following equivariant homotopy equivalences
\begin{gather*}
U \simeq \tilde{e}^0_0 = \pt, \qquad
V \simeq \tilde{e}^0_1 = D_6/D_2, \qquad
U \cap V \simeq \tilde{e}^1_{01} \simeq D_6/\Z_2^{(1)},
\end{gather*}
where $D_2 = \{ 1, C^3, \sigma_1, \sigma_4 \}$ and $\Z_2^{(1)} = \{ 1, \sigma_1 \}$. The equivariant cohomology groups of these spaces can be summarized as follows:
$$
\begin{array}{|c|c|c|c|}
\hline
n = 3 & & \Z_2 \oplus \Z_2 & 0 \\
\hline
n = 2 & & \Z_2^{\oplus 2} \oplus \Z_2^{\oplus 2} & \Z_2 \tsep{2pt}\bsep{2pt}\\
\hline
n = 1 & & 0 & 0 \\
\hline
n = 0 & & \Z \oplus \Z & \Z \\
\hline
& H^n_{D_6}(Y) & H^n_{D_6}(U) \oplus H^n_{D_6}(V) & H^n_{D_6}(U \cap V)\tsep{2pt}\bsep{2pt}\\
\hline
\end{array}
$$
In the Mayer--Vietoris exact sequence
\begin{gather*}
\cdots \to H^n_{D_6}(Y) \to H^n_{D_6}(U) \oplus H^n_{D_6}(V) \overset{\Delta}{\to} H^n_{D_6}(U \cap V) \to H^{n+1}_{D_6}(Y) \to \cdots,
\end{gather*}
the homomorphism $\Delta$ is expressed as $\Delta (u, v) = j_U^*(u) - j_V^*(v)$ with $j_U \colon U \cap V \to U$ and $j_V \colon U \cap V \to V$ the inclusions. This immediately determines $H^0_{D_6}(Y) \cong \Z$ and $H^1_{D_6}(Y) = 0$. To complete the proof, we recall the identif\/ications
\begin{gather*}
H^2_{D_6}(U) \cong \operatorname{Hom}(D_6, U(1)) \cong \Z_2^{\oplus 2}, \qquad
H^2_{D_6}(V) \cong \operatorname{Hom}(D_2, U(1)) \cong \Z_2^{\oplus 2}, \\
H^2_{D_6}(U \cap V) \cong \operatorname{Hom}\big(\Z_2^{(1)}, U(1)\big) \cong \Z_2,
\end{gather*}
under which $j_U^*$ and $j_V^*$ are induced from the inclusions $D_2 \to D_6$ and $\Z_2^{(1)} \to D_6$. As a basis of~$H^2_{D_6}(U)$ we can choose the following $1$-dimensional representations $\rho_i\colon D_6 \to U(1)$ of $D_6$
\begin{gather*}
\rho_1 \colon \
\begin{cases}
C \mapsto 1, \\
\sigma_1 \mapsto -1,
\end{cases}
\qquad
\rho_2 \colon \
\begin{cases}
C \mapsto -1, \\
\sigma_1 \mapsto 1.
\end{cases}
\end{gather*}
Similarly, we can choose the following $1$-dimensional representations $\rho'_i$ of $D_2 = \{ 1, \sigma_1, C^3, \sigma_4 \}$ as a basis of $H^2_{D_6}(V)$
\begin{gather*}
\rho'_1 \colon \
\begin{cases}
C^3 \mapsto 1, \\
\sigma_1 \mapsto -1,
\end{cases}
\qquad
\rho'_2 \colon \
\begin{cases}
C^3 \mapsto -1, \\
\sigma_1 \mapsto 1.
\end{cases}
\end{gather*}
Now, we can see $H^2_{D_6}(Y) \cong \operatorname{Ker}\Delta \cong \Z_2^{\oplus 3}$, and it has the following basis
\begin{gather*}
\{ (\rho_1, \rho'_1), (\rho_2, \rho_2'), (0, \rho'_2) \} \subset \operatorname{Hom}(D_6, U(1)) \oplus \operatorname{Hom}(D_4, U(1)).
\end{gather*}
We can also see that $\Delta$ is surjective, and $H^3_{D_6}(Y) \cong \Z_2^{\oplus 2}$.
\end{proof}

Let $X_1$ be the $1$-skeleton of the $D_6$-CW complex $T^2$.

\begin{lem}
$H^3_{D_6}(X_1; \Z) \cong \Z_2^{\oplus 2}$.
\end{lem}

\begin{proof}
We cover $X_1 = \tilde{e}^0_0 \cup \tilde{e}^0_1 \cup \tilde{e}^0_2 \cup \tilde{e}^1_{01} \cup \tilde{e}^1_{02} \cup \tilde{e}^1_{12}$ by invariant subspaces $U'$ and $V'$ which admit the following equivariant homotopy equivalences
\begin{gather*}
U' \simeq Y, \qquad
V' \simeq \tilde{e}^0_2 = D_6/D_3, \qquad
U' \cap V' \simeq \tilde{e}^1_{02} \sqcup \tilde{e}^1_{12} \simeq D_6/\Z_2^{(2)} \sqcup D_6/\Z_2^{(4)},
\end{gather*}
where $D_3 = \{ 1, C^2, C^4, \sigma_2, \sigma_4, \sigma_6 \}$, $\Z_2^{(2)} = \{ 1, \sigma_2 \}$ and $\Z_2^{(4)} = \{ 1, \sigma_4 \}$. The equivariant cohomology groups of these spaces are summarized as follows:
$$
\begin{array}{|c|c|c|c|}
\hline
n = 3 & & \Z_2^{\oplus 2} \oplus 0 & 0\tsep{2pt}\bsep{2pt} \\
\hline
n = 2 & & \Z_2^{\oplus 3} \oplus \Z_2 & \Z_2 \oplus \Z_2 \tsep{2pt}\bsep{2pt}\\
\hline
n = 1 & & 0 & 0 \\
\hline
n = 0 & & \Z \oplus \Z & \Z \oplus \Z \\
\hline
& H^n_{D_6}(X_1) & H^n_{D_6}(U') \oplus H^n_{D_6}(V') & H^n_{D_6}(U' \cap V') \tsep{2pt}\bsep{2pt}\\
\hline
\end{array}
$$
The homomorphism $\Delta$ in the Mayer--Vietoris exact sequence
\begin{gather*}
\cdots \to H^2_{D_6}(X_1) \to H^2_{D_6}(U') \oplus H^2_{D_6}(V') \overset{\Delta}{\to} H^2_{D_6}(U' \cap V') \to H^3_{D_6}(X_1) \to \cdots
\end{gather*}
is expressed as $\Delta(u, v) = j_{U'}^*(u) - j_{V'}^*(v)$ by using the inclusions $j_{U'} \colon U' \cap V' \to U'$ and $j_{V'} \colon U' \cap V' \to V'$. An inspection proves that $j_{U'}^*$ agrees with the composition of the following two homomorphisms:
\begin{itemize}\itemsep=0pt
\item[(i)]
the inclusion that follows from the calculation of $H^2_{D_6}(Y)$ in Lemma \ref{lem:p6m_Y}
\begin{gather*}
H^2_{D_6}(U') \cong H^2_{D_6}(Y) \longrightarrow
\operatorname{Hom}(D_6, U(1)) \oplus \operatorname{Hom}(D_2, U(1)).
\end{gather*}

\item[(ii)]
the direct sum $i_2^* \oplus i_4^*$ of the homomorphisms
\begin{gather*}
i^*_2 \colon \ \operatorname{Hom}(D_6, U(1)) \to \operatorname{Hom}\big(\Z_2^{(2)}, U(1)\big), \qquad
i^*_4 \colon \ \operatorname{Hom}(D_2, U(1)) \to \operatorname{Hom}\big(\Z_2^{(4)}, U(1)\big),
\end{gather*}
induced from the inclusions $i_2\colon \Z_2^{(2)} \to D_6$ and $\Z_2^{(4)} \to D_2$.

\end{itemize}
Then, using the basis presented in the calculation of $H^2_{D_6}(Y)$, we f\/ind
\begin{gather*}
j_{U'}^*(\rho_1, \rho_1') = (\rho, \rho), \qquad j_{U'}^*(\rho_2, \rho_2') = (\rho, \rho), \qquad
j_{U'}^*(0, \rho_2') = (0, \rho),
\end{gather*}
where $\rho\colon \Z_2 \to \Z_2$ is the identity map generating $\operatorname{Hom}(\Z_2, U(1)) \cong \Z_2$. Hence $j_{U'}^*$ as well as $\Delta$ are surjective, and $H^3_{D_6}(X_1) \cong H^3_{D_6}(Y) \cong \Z_2^{\oplus 2}$.
\end{proof}

\begin{thm}[\textsf{p6m}]
$H^3_{D_6}(T^2; \Z) \cong \Z_2^{\oplus 2}$.
\end{thm}

\begin{proof}
The relevant part of the exact sequence for the pair $(T^2, X_1)$ is
\begin{gather*}
H^3_{D_6}\big(T^2, X_1\big) \to H^3_{D_6}\big(T^2\big) \to H^3_{D_6}(X_1) \to H^4_{D_6}\big(T^2, X_1\big).
\end{gather*}
By means of the excision axiom, we have $H^n_{D_6}(T^2, X_1) \cong H^{n-2}(\pt)$. Therefore we get $H^3_{D_6}(T^2) \cong H^3_{D_6}(X_1) \cong \Z_2^{\oplus 2}$.
\end{proof}

Let $\hat{\Z} = \Z_{\phi_1}$ be the $D_6$-module such that its underlying group is $\Z$ and $D_6$ acts via the homomorphism $\phi_1\colon D_6 \to \Z_2$ given by $\phi_1(C) = -1$ and $\phi_1(\sigma_1) = 1$.

\begin{lem} \label{lem:p6m_short_exact_sequence_of_modules}
There is an exact sequence of $D_6$-modules
\begin{gather*}
0 \to H^1\big(T^2; \Z\big) \to H^1(Y; \Z) \overset{\pi}{\to} \hat{\Z} \to 0
\end{gather*}
admitting a module homomorphism $s\colon \hat{\Z} \to H^1(Y; \Z)$ such that $\pi \circ s = 3$.
\end{lem}

\begin{proof}
Let $\eta_1, \eta_2 \in H_1(T^2)$ be the homology classes of the loops going along the vectors $1$ and $\tau$ respectively in the fundamental domain, which form a basis of $H_1(T^2) \cong \Z^2$. Also, let $\gamma_1, \gamma_2, \gamma_3 \in H_1(Y)$ be the homology classes of loops along $1$, $\tau$ and $\tau - 1$, which form a basis of $H_1(Y) \cong \Z^3$. The inclusion map $i\colon Y \to T^2$ relates these bases by $i_*\gamma_1 = \eta_1$, $i_*\gamma_2 = \eta_2$ and $i_*\gamma_3 = \eta_2 - \eta_1$. The actions of $C \in D_6$ and $\sigma_1$ on these bases are
\begin{gather*}
\begin{cases}
C_*\eta_1 = \eta_2, \\
C_*\eta_2 = \eta_2 - \eta_1,
\end{cases}
\qquad
\begin{cases}
{\sigma_1}_*\eta_1 = \eta_1, \\
{\sigma_1}_*\eta_2 = \eta_1 - \eta_2,
\end{cases}
\qquad
\begin{cases}
C_*\gamma_1 = \gamma_2, \\
C_*\gamma_2 = \gamma_3, \\
C_*\gamma_3 = - \gamma_1,
\end{cases}
\qquad
\begin{cases}
{\sigma_1}_*\gamma_1 = \gamma_1, \\
{\sigma_1}_*\gamma_2 = - \gamma_3, \\
{\sigma_1}_*\gamma_3 = - \gamma_2.
\end{cases}\!\!\!
\end{gather*}
Let $\{ h_1, h_2 \} \subset H^1(T^2)$ and $\{ g_1, g_2, g_3 \} \subset H^1(Y)$ be dual to the homology bases. They are related by $i^*h_1 = g_1 - g_3$ and $i^*h_2 = g_2 + g_3$, and the induced $D_6$-actions are as follows.
\begin{gather*}
\begin{cases}
C^*h_1 = - h_2, \\
C^*h_2 = h_1 + h_2,
\end{cases}\qquad
\begin{cases}
\sigma_1^*h_1 = h_1 + h_2, \\
\sigma_1^*h_2 = - h_2,
\end{cases}\qquad
\begin{cases}
C^*g_1 = - g_3, \\
C^*g_2 = g_1, \\
C^*g_3 = g_2,
\end{cases}\qquad
\begin{cases}
\sigma_1^*g_1 = g_1, \\
\sigma_1^*g_2 = - g_3, \\
\sigma_1^*g_3 = - g_2.
\end{cases}
\end{gather*}
These expressions allow us to prove that the cokernel of the homomorphism $i^* \colon H^1(T^2) \to H^2(Y)$ is isomorphic to $\hat{\Z}$, yielding the exact sequence. The homomorphism $s\colon \hat{\Z} \to H^1(Y)$ is given by $s(1) = g_1 - g_2 + g_3$.
\end{proof}

\begin{lem}
$H^n_{\mathrm{group}}(D_6; H^1(T^2; \Z)) = 0$ for $n = 0, 1, 2$.
\end{lem}

\begin{proof}
We use the long exact sequence in group cohomology induced from the exact sequence $0 \to H^1(T^2) \to H^1(Y) \overset{\pi}{\to} \hat{\Z} \to 0$ in coef\/f\/icients. By Lemmas~\ref{lem:1_dim_complex}, \ref{lem:p6m_Y} and~\ref{lem:cohomology_of_point_twisted_case} to be given in Section~\ref{sec:twisted_case}, we get the following:
$$
\begin{array}{|c|c|c|c|}
\hline
n = 2 & & \Z_2 & \Z_2 \\
\hline
n = 1 & & \Z_2 & \Z_2 \\
\hline
n = 0 & & 0 & 0 \\
\hline
& H^n_{\mathrm{group}}(D_6; H^1(T^2)) & H^n_{\mathrm{group}}(D_6; H^1(Y)) &
H^n_{\mathrm{group}}(D_6; \hat{\Z})\tsep{2pt}\bsep{2pt}\\
\hline
\end{array}
$$
It is clear that $H^0_{\mathrm{group}}(D_6; H^1(T^2)) = 0$. The homomorphism in group cohomology induced from $\pi\colon H^1(Y) \to \hat{\Z}$ is surjective in degree $1$ and $2$, because $\pi \circ s = 3$. This leads to the remaining vanishing.
\end{proof}

\begin{thm}[\textsf{p6m}]
The following holds true:
\begin{itemize}\itemsep=0pt
\item[$(a)$] $F^2H^3_{D_6}(T^2; \Z) \cong \Z_2$,

\item[$(b)$] the $D_6$-equivariant cohomology of $T^2$ in low degrees is as follows:
\begin{gather*}
H^0_{D_6}\big(T^2; \Z\big) \cong \Z, \qquad H^1_{D_6}\big(T^2; \Z\big) = 0, \qquad H^2_{D_6}\big(T^2; \Z\big) \cong \Z_2^{\oplus 2}.
\end{gather*}
\end{itemize}
\end{thm}

\begin{proof} In the $E_2$-term of the Leray--Serre spectral sequence $E_2^{p, q} = H^p_{\mathrm{group}}(D_6; H^q(T^2; \Z))$, the coef\/f\/icient $H^0(T^2)$ is identif\/ied with the trivial $D_6$-module $\Z$, and $H^2(T^2)$ with $\tilde{\Z}$ as in Lemma~\ref{lem:group_coh_orientation_coeff}. The group cohomology with coef\/f\/icients in $H^1(T^2)$ is already computed, and that in $\tilde{\Z}$ is also computed in Lemma~\ref{lem:group_coh_orientation_coeff}. The $E_2$-terms are summarized as follows:
$$
\begin{array}{|c|c|c|c|c|}
\hline
q = 3 & 0 & 0 & 0 & 0 \\
\hline
q = 2 & 0 & \Z_2 & & \\
\hline
q = 1 & 0 & 0 & 0 & \\
\hline
q = 0 & \Z & 0 & \Z_2^{\oplus 2} & \Z_2 \tsep{2pt}\bsep{2pt}\\
\hline
E_2^{p, q} & p = 0 & p = 1 & p = 2 & p = 3 \tsep{2pt}\bsep{2pt}\\
\hline
\end{array}
$$
This list and Lemma \ref{lem:typical_argument_non_orientation_preserving_case} lead to the theorem.
\end{proof}
